\chapter*{Abstract}
\addcontentsline{toc}{chapter}{Abstract}

Room rates on Online Travel Agencies change as inventory, lead time and market conditions evolve, yet a single price snapshot cannot show whether a listed rate is rising or falling. This proposal describes a web-based system that repeatedly collects publicly visible Booking.com room rates for a selected cohort of 355 properties in five Vietnamese cities, forecasts short-term rate movements with machine learning, and serves the forecasts through two decision layers: a hotelier view centred on competitive-set monitoring and pricing-position alerts, and a consumer view providing informational booking-timing guidance. The system does not process bookings or guarantee a lowest price.

One modelling requirement shapes the entire data pipeline. Every price record carries two independent time coordinates, the observation timestamp and the stay date, so the dataset records how the rate for a fixed check-in date evolves as that date approaches. Lead time is derived from the two coordinates and becomes a central predictor. A durable MySQL work queue drives a Selenium crawler that separates sold-out nights from properties that cannot be booked at all and from dead listings. A reference definition calibrated per hotel and per check-in date keeps consecutive observations describing a comparable room and rate plan. From 18 August to 30 November 2026, the operator follows a fixed calendar and uses the scraper web interface to start one full-cohort run per day with nine near-term and three farther check-in dates, providing dense short-lead series while retaining longer-horizon coverage.

The required modelling deliverable is a reproducible regression pipeline for the 1, 3, 7 and 14-day horizons. Persistence and simple statistical baselines are compared with Random Forest and XGBoost; the tree models are tuned with time-aware cross-validation using RandomizedSearchCV followed by a focused GridSearchCV. Deep-learning and standalone classification models are optional extensions only if the core regression system is complete. Evaluation uses a strictly chronological split and adopts \emph{Accuracy@20\%} as the primary feasibility measure: at least 80 per cent of held-out predictions should have an absolute percentage error no greater than 20 per cent. MAE, RMSE, sMAPE, MAPE, $R^2$, directional accuracy, SHAP attribution and two retrospective decision simulations provide complementary evidence.

\clearpage
